\section*{Bab \ref{ch:b2:multint} --- Integral Garis dan Integral Lipat}

\begin{solution}{exo:b2:multint:1}
(a) Ruas $\gamma(t) = (t, t)$ dengan $t \in [0,1]$:
\[
\int_\gamma y^2\dd x + x\dd y
= \int_0^1 (t^2 + t)\,\dd t = \frac13 + \frac12 = \frac56 .
\]
(b) Parabola $\gamma(t) = (t, t^2)$:
\[
\int_0^1 \bigl(t^4\cdot 1 + t\cdot 2t\bigr)\dd t
= \frac15 + \frac23 = \frac{13}{15} .
\]
Kedua nilainya berbeda, jadi integralnya bergantung pada lintasannya: sehingga bentuknya
\emph{tak} \omterm{def:b2:multint:exact}{eksak} --- dan secara konsisten, $P_y = 2y \neq 1 = Q_x$, jadi
ia bahkan tak tertutup.
\end{solution}

\begin{solution}{exo:b2:multint:2}
Berlaku $P = 2xy + y^3$ dan $Q = x^2 + 3xy^2 + 1$: jadi $P_y = 2x + 3y^2 = Q_x$,
sehingga tertutup pada $\R^2$. Carilah $f$ dengan $f_x = P$: maka $f = x^2y + xy^3
+ g(y)$; lalu $f_y = x^2 + 3xy^2 + g'(y) = Q$ memaksa $g'(y) = 1$,
katakanlah $g(y) = y$. Jadi
\[
f(x, y) = x^2y + xy^3 + y
\]
merupakan sebuah \omterm{def:b2:multint:exact}{potensial} (sebab $\R^2$ berbentuk bintang, jadi sebuah \omterm{def:b2:multint:exact}{potensialnya} mesti
ada menurut lema Poincaré --- tetapi memamerkannya lebih cepat). Menurut
\cref{thm:b2:multint:ftc}, untuk sembarang busur dari $(0,0)$ ke $(1,2)$:
\[
\int_\gamma\omega = f(1, 2) - f(0, 0) = 2 + 8 + 2 = 12 .
\]
\end{solution}

\begin{solution}{exo:b2:multint:3}
Daerahnya adalah $\{0 \leq x \leq 1,\ x^2 \leq y \leq x\}$.
Dengan $y$ dulu:
\[
\int_0^1\!\int_{x^2}^{x}(x + y)\,\dd y\,\dd x
= \int_0^1\Bigl(x(x - x^2)
+ \frac{x^2 - x^4}{2}\Bigr)\dd x
= \int_0^1\Bigl(\frac{3x^2}{2} - x^3 - \frac{x^4}{2}\Bigr)\dd x
= \frac12 - \frac14 - \frac1{10} = \frac{3}{20} .
\]
Dengan $x$ dulu: irisannya pada ketinggian $y \in [0, 1]$ adalah $y \leq x \leq
\sqrt y$, jadi
\[
\int_0^1\!\int_{y}^{\sqrt y}(x + y)\,\dd x\,\dd y
= \int_0^1\Bigl(\frac{y - y^2}{2}
+ y(\sqrt y - y)\Bigr)\dd y
= \frac14 - \frac16 + \frac25 - \frac13 = \frac{3}{20} .
\]
\end{solution}

\begin{solution}{exo:b2:multint:4}
\emph{Integral pertamanya.} Pada cakram $D_R$, dalam koordinat kutub:
\[
\iint_{D_R}\frac{\dd x\,\dd y}{(1 + x^2 + y^2)^2}
= \int_0^{2\pi}\!\!\int_0^R
\frac{\rho\,\dd\rho\,\dd\alpha}{(1 + \rho^2)^2}
= 2\pi\Bigl[-\frac{1}{2(1 + \rho^2)}\Bigr]_0^R
= \pi\Bigl(1 - \frac{1}{1 + R^2}\Bigr)
\xrightarrow[R \to \infty]{} \pi .
\]
\emph{Integral keduanya.} Seperempat cakramnya adalah $0 \leq \alpha \leq
\frac\pi2$ dan $0 \leq \rho \leq 1$, dengan $xy =
\rho^2\cos\alpha\sin\alpha$:
\[
\iint_{D'}xy\,\dd x\,\dd y
= \int_0^{\pi/2}\!\!\cos\alpha\sin\alpha\,\dd\alpha
\int_0^1 \rho^3\,\dd\rho
= \frac12\cdot\frac14 = \frac18 .
\]
\end{solution}

\begin{solution}{exo:b2:multint:5}
Menurut \cref{cor:b2:multint:area} dengan $x = \cos^3 t$ dan $y = \sin^3 t$:
berlaku $x' = -3\cos^2 t\sin t$ dan $y' = 3\sin^2 t\cos t$, jadi
\[
xy' - yx' = 3\cos^4 t\sin^2 t + 3\sin^4 t\cos^2 t
= 3\sin^2 t\cos^2 t = \frac{3}{4}\sin^2 2t
= \frac{3}{8}(1 - \cos 4t) .
\]
Karena itu
\[
A = \frac12\int_0^{2\pi}\frac38(1 - \cos 4t)\,\dd t
= \frac{3}{16}\cdot 2\pi = \frac{3\pi}{8} .
\]
(\omterm{pb:b2:curves:1}{Astroidnya} muat di cakram satuan yang berluas $\pi$; jadi tiga per delapan
$\pi$ masuk akal bagi bentuk bintangnya yang beruncing empat.)
\end{solution}

\begin{solution}{exo:b2:multint:6}
\emph{Penumpukannya:} di atas tiap $(x, y)$ pada cakram satuan $D$, peubah $z$
berjalan dari $x^2 + y^2$ ke $1$:
\[
V = \iint_D \bigl(1 - x^2 - y^2\bigr)\dd x\,\dd y
= \int_0^{2\pi}\!\!\int_0^1 (1 - \rho^2)\rho\,\dd\rho\,\dd\alpha
= 2\pi\Bigl(\frac12 - \frac14\Bigr) = \frac\pi2 .
\]
\emph{Pengirisannya:} irisannya pada ketinggian $z \in [0, 1]$ adalah cakram
$x^2 + y^2 \leq z$ yang berluas $\pi z$:
\[
V = \int_0^1 \pi z\,\dd z = \frac\pi2 .
\]
\end{solution}

\begin{solution}{exo:b2:multint:7}
Dalam koordinat bola unsur volumenya adalah
$r^2\cos\varphi\,\dd r\,\dd\theta\,\dd\varphi$
(\cref{ex:b2:multint:spherical}), dan bagian sudutnya berintegral
$4\pi$ (yakni $2\pi$ dari $\theta$, dan $\int_{-\pi/2}^{\pi/2}\cos = 2$).
Jadi volume cangkangnya adalah
\[
\int_a^b 4\pi r^2\,\dd r = \frac43\pi\bigl(b^3 - a^3\bigr).
\]
Adapun untuk integral keduanya, integran $1/r$ hanya bergantung pada $r$:
\[
\iiint_B \frac{\dd x\,\dd y\,\dd z}{r}
= \int_0^R 4\pi r^2\cdot\frac1r\,\dd r
= 4\pi\,\frac{R^2}{2} = 2\pi R^2 .
\]
(Integrannya meledak di titik asalnya, tetapi tanpa bahaya: sebab $r^2/r = r$
bersifat \omterm{def:b2:metric:continuity}{kontinu} --- jadi integral atas cangkang $\varepsilon \leq r
\leq R$ konvergen ketika $\varepsilon \to 0$, dan itulah makna
persisnya pernyataan itu. Adapun perhitungan semacam ini adalah langkah
pertama menuju teorema Newton bahwa bola homogen menarik seperti
massa titik di pusatnya.)
\end{solution}

\begin{solution}{exo:b2:multint:8}
Integran $g(t; x, y) = xP(tx, ty) + yQ(tx, ty)$ bersifat
$\mathcal{C}^1$ dalam $(x, y)$, \omterm{def:b2:metric:continuity}{kontinu} dalam $t$, dengan turunan
parsialnya \omterm{def:b2:metric:continuity}{kontinu} pada $[0,1] \times U$; jadi penurunan di bawah
tanda integral (\cref{ch:b2:integration}, yang diterapkan pada
interval-$t$ \omterm{def:b2:metric:compact}{kompak} $[0,1]$, dengan dominasinya otomatis di sana)
memberi
\[
f_x(x, y)
= \int_0^1 \bigl(P(tx, ty) + tx\,P_x(tx, ty)
+ ty\,Q_x(tx, ty)\bigr)\dd t .
\]
Lalu dengan memakai ketertutupannya $Q_x = P_y$:
\[
tx\,P_x(tx, ty) + ty\,P_y(tx, ty)
= t\,\frac{\dd}{\dd t}\bigl[P(tx, ty)\bigr] ,
\]
jadi integrannya adalah $P(tx, ty) + t\frac{\dd}{\dd t}P(tx, ty) =
\frac{\dd}{\dd t}\bigl[t\,P(tx, ty)\bigr]$ dan
\[
f_x(x, y) = \Bigl[t\,P(tx, ty)\Bigr]_0^1 = P(x, y) .
\]
Secara simetris $f_y = Q$ (lewat perhitungan yang sama dengan $P_y = Q_x$ dipakai
pada arah yang lain). Perhatikanlah di mana hipotesisnya masuk: sebab $f$ didefinisikan
dengan mengintegralkan sepanjang ruas $[0, M]$, yang terletak di $U$
persis karena $U$ berbentuk bintang.
\end{solution}

\begin{solution}{exo:b2:multint:9}
Pada jalur $[0, A] \times [0, \infty)$ fungsi $(x, y)
\mapsto e^{-xy}\sin x$ tak terintegralkan secara \omterm{def:b2:series:def}{mutlak} sampai $y =
\infty$ secara seragam dalam pengertian yang naif, tetapi tiap integral teriterasinya
konvergen dan kesamaannya menyusul dari Fubini pada $[0, A]
\times [0, B]$ ditambah sebuah limit $B \to \infty$ (sebab ekornya
$\int_0^A\int_B^\infty e^{-xy}\abs{\sin x}\,\dd y\,\dd x \leq
\int_0^A \frac{e^{-Bx}\abs{\sin x}}{x}\dd x \leq \int_0^A e^{-Bx}
\dd x\to 0$, dengan memakai $\abs{\sin x} \leq x$).

\emph{Dengan $y$ dulu:} berlaku $\int_0^\infty e^{-xy}\,\dd y = \frac1x$ untuk $x
> 0$, jadi integral pertamanya adalah $\int_0^A \frac{\sin x}{x}\,\dd x$.

\emph{Dengan $x$ dulu:} dua kali pengintegralan parsial (atau mengambil
bagian imajiner $\int_0^A e^{(i - y)x}\dd x$) memberi
\[
\int_0^A e^{-xy}\sin x\,\dd x
= \frac{1 - e^{-Ay}(y\sin A + \cos A)}{1 + y^2} .
\]
Lalu setelah diintegralkan dalam $y$ atas $[0, \infty)$, suku $\int_0^\infty
\frac{\dd y}{1 + y^2} = \frac\pi2$ terpisah:
\[
\int_0^A \frac{\sin x}{x}\,\dd x
= \frac{\pi}{2}
- \int_0^\infty e^{-Ay}\,\frac{y\sin A + \cos A}{1 + y^2}\,\dd y .
\]
Adapun sisanya dibatasi oleh $\int_0^\infty e^{-Ay}\frac{y +
1}{1 + y^2}\dd y \leq \int_0^\infty e^{-Ay}\cdot\frac{1+y}{1+y^2}
\,\dd y \to 0$ ketika $A \to \infty$ (lewat kekonvergenan terdominasi, atau
lewat batas kasar $\frac{1 + y}{1 + y^2} \leq \frac32$ yang memberi
$\frac{3}{2A}$). Karena itu $\int_0^\infty\frac{\sin x}{x}\dd x =
\frac\pi2$ --- yakni nilai yang sama yang diperoleh pada
\cref{ch:b2:integration} lewat menurunkan \omterm{thm:b2:integration:continuity}{integral berparameter};
sedangkan di sini Fubini yang mengerjakannya.
\end{solution}

\begin{solution}{exo:b2:multint:10}
Parameterkan lewat \omterm{def:b2:curves:length}{panjang busurnya} $s \in [0, 2\pi]$, jadi $x'^2 + y'^2 = 1$,
lalu geserlah sehingga $\int_0^{2\pi} x(s)\,\dd s = 0$. Menurut
\cref{cor:b2:multint:area},
\[
2A = \oint x\,\dd y - y\,\dd x
= \int_0^{2\pi}\bigl(xy' - yx'\bigr)\dd s .
\]
Lalu mengintegralkan $\oint y\,\dd x$ secara parsial atas kalanya (dengan suku
perbatasannya saling meniadakan berkat kekalaannya), $-\int yx' = \int y'x$, jadi sebenarnya
$2A = 2\int_0^{2\pi}xy'\,\dd s$. Kemudian $2xy' \leq x^2 + y'^2$
memberi
\[
2A \leq \int_0^{2\pi}\bigl(x^2 + y'^2\bigr)\dd s
= \int_0^{2\pi} x^2 + \int_0^{2\pi}\bigl(1 - x'^2\bigr)
= 2\pi - \int_0^{2\pi}\bigl(x'^2 - x^2\bigr)\dd s .
\]
Ketaksamaan Wirtinger (\cref{ch:b2:fourier}, pada latihannya: yakni untuk
fungsi berkala-$2\pi$ berkelas $\mathcal{C}^1$ yang bererata nol,
$\int x^2 \leq \int x'^2$) membuat integral terakhirnya taknegatif:
jadi $A \leq \pi$. Adapun kesamaannya menuntut kesamaan pada Wirtinger (yakni $x(s) =
a\cos s + b\sin s$) dan pada $2xy' \leq x^2 + y'^2$ (yakni $y' = x$
\omterm{def:b2:funcseq:def}{titik demi titik}), yang memaksa $y = a\sin s - b\cos s + c$: jadi kurvanya adalah
lingkaran satuan (yang dipusatkan sepantasnya). Lalu karena kurva berpanjang $L$
menskala ke \omterm{def:b2:curves:length}{panjang} $2\pi$, pernyataan umumnya adalah $A \leq
\frac{L^2}{4\pi}$: jadi di antara semua kurva tertutup berkeliling tertentu,
lingkarannya melingkupi \omterm{def:b2:multint:domain}{luas} terbesar.
\end{solution}

\begin{solution}{exo:b2:multint:11}
Dalam koordinat bola berlaku $z = r\sin\varphi$ dan $\dd x\,\dd
y\,\dd z = r^2\cos\varphi\,\dd r\,\dd\theta\,\dd\varphi$:
\[
\iiint_B z^2
= \int_0^R r^4\,\dd r\int_0^{2\pi}\dd\theta
\int_{-\pi/2}^{\pi/2}\sin^2\varphi\cos\varphi\,\dd\varphi
= \frac{R^5}5\cdot2\pi\cdot
\Bigl[\frac{\sin^3\varphi}3\Bigr]_{-\pi/2}^{\pi/2}
= \frac{4\pi R^5}{15}.
\]
Berkat kesimetrian bolanya di bawah permutasi koordinatnya,
$\iiint_B x^2 = \iiint_B y^2 = \iiint_B z^2$, jadi $\iiint_B(x^2
+ y^2 + z^2) = 3\cdot\frac{4\pi R^5}{15} = \frac{4\pi R^5}5$.
Pemeriksaan cangkangnya: $\int_0^R r^2\cdot 4\pi r^2\,\dd r = \frac{4\pi
R^5}5$ --- sebab integran $r^2$ bersifat tetap pada bola berjari-jari
$r$ yang berluas $4\pi r^2$.
\end{solution}

\begin{solution}{exo:b2:multint:12}
Parameterkan rusuk dari $(x_i, y_i)$ ke $(x_{i+1}, y_{i+1})$
lewat $\gamma(t) = \bigl((1-t)x_i + tx_{i+1},\ (1-t)y_i +
ty_{i+1}\bigr)$. Maka sumbangannya pada
$\frac12\oint(x\,\dd y - y\,\dd x)$ adalah
\[
\frac12\int_0^1\Bigl(\bigl((1-t)x_i +
tx_{i+1}\bigr)(y_{i+1} - y_i) - \bigl((1-t)y_i +
ty_{i+1}\bigr)(x_{i+1} - x_i)\Bigr)\dd t ,
\]
lalu karena $\int_0^1\bigl((1-t)u + tv\bigr)\dd t = \frac{u +
v}2$, ini sama dengan
\[
\frac14\Bigl((x_i + x_{i+1})(y_{i+1} - y_i) - (y_i +
y_{i+1})(x_{i+1} - x_i)\Bigr)
= \frac12\bigl(x_iy_{i+1} - x_{i+1}y_i\bigr),
\]
sebab suku silangnya saling meniadakan. Lalu menjumlahkannya atas $m$ rusuknya memberi
rumus tali sepatunya, menurut \cref{cor:b2:multint:area}. Adapun segitiga
$(0,0), (1,0), (0,1)$: $\frac12\bigl((0\cdot0 - 1\cdot0) +
(1\cdot1 - 0\cdot0) + (0\cdot0 - 0\cdot1)\bigr) = \frac12$,
yakni luasnya yang benar.
\end{solution}

\begin{solution}{pb:b2:multint:1}
\textbf{1.} Bola $B_n(R)$ diperikan lewat batas teriterasi
$-R \leq x_n \leq R$, lalu $\abs{x_{n-1}} \leq \sqrt{R^2 -
x_n^2}$, dan seterusnya; jadi menyulihkan $x_i = Ru_i$ pada masing-masing
$n$ integral satu peubahnya mengalikan tiap-tiapnya dengan $R$ lalu memetakan
batasnya ke batas $B_n(1)$: sehingga $v_n(R) = R^n\,v_n(1) =
V_nR^n$.

\textbf{2.} Dengan mengiris sepanjang $x_n = t$: irisan $B_n(1)$ adalah
bola $B_{n-1}\bigl(\sqrt{1 - t^2}\bigr)$, jadi menurut pertanyaan
1,
\[
V_n = \int_{-1}^1 v_{n-1}\bigl(\sqrt{1 - t^2}\bigr)\,\dd t
= V_{n-1}\int_{-1}^{1}(1 - t^2)^{\frac{n-1}2}\,\dd t .
\]

\textbf{3.} Dengan $t = \sin\theta$: berlaku $\dd t =
\cos\theta\,\dd\theta$ dan $(1 - t^2)^{\frac{n-1}2} =
\cos^{n-1}\theta$ pada $\intcc{-\pi/2}{\pi/2}$, jadi
\[
\int_{-1}^1(1 - t^2)^{\frac{n-1}2}\dd t
= \int_{-\pi/2}^{\pi/2}\cos^n\theta\,\dd\theta
= 2\int_0^{\pi/2}\cos^n\theta\,\dd\theta = 2W_n,
\]
lalu $\theta \mapsto \frac\pi2 - \theta$ mempertukarkan bentuk sinus
dan kosinus milik $W_n$.

\textbf{4.} Tulislah $\sin^n = \sin^{n-2}(1 - \cos^2)$ lalu
integralkan $\int\sin^{n-2}\cos\cdot\cos$ secara parsial (dengan $v =
\frac{\sin^{n-1}}{n-1}$):
\[
W_n = W_{n-2} - \frac{W_n}{n-1}
\quad\Longrightarrow\quad
W_n = \frac{n-1}{n}W_{n-2}.
\]
Karena itu $nW_nW_{n-1} = (n-1)W_{n-1}W_{n-2}$: jadi barisan
$(nW_nW_{n-1})$ bersifat tetap dan sama dengan $1\cdot W_1W_0 =
1\cdot\frac\pi2$, sehingga $W_nW_{n-1} = \frac{\pi}{2n}$.

\textbf{5.} Pertanyaan 2--3 memberi $V_n = 2W_nV_{n-1}$, dua kali:
\[
V_n = 2W_n\cdot 2W_{n-1}\,V_{n-2}
= 4\,\frac{\pi}{2n}\,V_{n-2} = \frac{2\pi}n\,V_{n-2}.
\]

\textbf{6.} Dari $V_0 = 1$: berlaku $V_{2k} = \frac{2\pi}{2k}V_{2k-2}
= \frac\pi kV_{2k-2}$, jadi $V_{2k} = \frac{\pi^k}{k!}$ secara
induktif. Sedangkan dari $V_1 = 2$: $V_{2k+1} =
\frac{2\pi}{2k+1}V_{2k-1}$, jadi
\[
V_{2k+1} = 2\prod_{j=1}^k\frac{2\pi}{2j+1}
= \frac{2^{k+1}\pi^k}{1\cdot3\cdot5\cdots(2k+1)} .
\]

\textbf{7.} Berlaku $V_1 = 2$, $V_2 = \pi \approx 3.142$, $V_3 =
\frac{4\pi}3 \approx 4.189$, $V_4 = \frac{\pi^2}2 \approx
4.935$, $V_5 = \frac{8\pi^2}{15} \approx 5.264$, $V_6 =
\frac{\pi^3}6 \approx 5.168$, dan $V_7 = \frac{16\pi^3}{105}
\approx 4.725$. Adapun nisbah satu langkahnya adalah $V_n/V_{n-1} = 2W_n$,
dan $(W_n)$ bersifat turun (sebab $\sin^n \leq \sin^{n-1}$
\omterm{def:b2:funcseq:def}{titik demi titik}). Kini $2W_5 = 2\cdot\frac45\cdot\frac23 =
\frac{16}{15} > 1$ sedangkan $2W_6 =
2\cdot\frac56\cdot\frac34\cdot\frac12\cdot\frac\pi2 =
\frac{5\pi}{16} < 1$: jadi nisbahnya melampaui $1$ sampai $n = 5$ lalu
berada di bawah $1$ sejak $n = 6$ --- sehingga $(V_n)$ naik ke
maksimumnya $V_5$ lalu turun.

\textbf{8.} Untuk $n \geq 13 > 4\pi$: berlaku $V_n/V_{n-2} = 2\pi/n <
\tfrac12$, jadi $V_{n} \leq C\cdot 2^{-n/2}$ dengan sebuah konstanta
yang tetap; bahkan lebih baik, untuk sembarang $q > 0$ berlaku $2\pi/n < q^2$ untuk $n$
yang besar, jadi $V_n/q^n \to 0$: sehingga peluruhannya mengalahkan setiap barisan
geometri. Adapun kekonvergenan $\sum V_n$ menyusul dari nisbah
$V_n/V_{n-2} \to 0$ (yakni bandingkan dengan deret geometri sejak
suatu pangkat).

\textbf{9.} Berlaku $\sum_{k\geq0}V_{2k}x^{2k} =
\sum_{k\geq0}\frac{(\pi x^2)^k}{k!} = \eu^{\pi x^2}$, yakni
deret eksponensialnya (\cref{ch:b2:powerseries}), yang konvergen untuk
setiap $x$. Lalu di $x = 1$: $\sum_kV_{2k} = \eu^\pi \approx
23.14$.

\textbf{10.} Pada kubus $\intcc{-R}R^n$ integrannya adalah
hasil kali $\prod_i\eu^{-x_i^2}$, jadi integral teriterasinya
terfaktorkan: yakni $\bigl(\int_{-R}^R\eu^{-t^2}\dd t\bigr)^n$.
Lalu dengan melewatkan $R \to \infty$ dan memakai
$\int_\R\eu^{-t^2}\dd t = \sqrt\pi$
(\cref{ex:b2:multint:polar}): jadi $I_n = \pi^{n/2}$.

\textbf{11.} Sulihan $t = u^2$ memberi $\Gamma(\tfrac12) =
\int_0^\infty t^{-1/2}\eu^{-t}\dd t =
2\int_0^\infty\eu^{-u^2}\dd u = \sqrt\pi$. Lalu dengan mengiterasikan
$\Gamma(s+1) = s\Gamma(s)$: berlaku $\Gamma(k+1) = k!\,\Gamma(1) =
k!$, dan
\[
\Gamma\Bigl(k + \frac32\Bigr)
= \Bigl(k + \frac12\Bigr)\Bigl(k - \frac12\Bigr)\cdots
\frac12\cdot\Gamma\Bigl(\frac12\Bigr)
= \frac{(2k+1)(2k-1)\cdots1}{2^{k+1}}\,\sqrt\pi .
\]

\textbf{12.} Tetapkan $F_n = \pi^{n/2}/\Gamma(\frac n2 + 1)$.
Karena $\Gamma(\frac n2 + 1) = \frac n2\,\Gamma(\frac n2) =
\frac n2\,\Gamma(\frac{n-2}2 + 1)$, kita memperoleh $F_n =
\frac{2\pi}nF_{n-2}$: yakni rekursi yang sama seperti $V_n$ (pertanyaan
5). Adapun basisnya: $F_1 = \sqrt\pi/\Gamma(\frac32) =
\sqrt\pi/(\frac{\sqrt\pi}2) = 2 = V_1$ dan $F_2 =
\pi/\Gamma(2) = \pi = V_2$. Jadi secara induktif $V_n = F_n$ untuk setiap
$n$; lalu pertanyaan 11 mengubahnya kembali menjadi kedua \omterm{def:b2:multint:exact}{bentuk tertutup}
pertanyaan 6.

\textbf{13.} Dengan $r = \sqrt t$: berlaku $\int_0^\infty
\eu^{-r^2}r^{n-1}\dd r = \frac12\int_0^\infty
t^{\frac n2 - 1}\eu^{-t}\dd t = \frac12\Gamma(\frac n2)$.
Karena itu
\[
n\,V_n\int_0^\infty\eu^{-r^2}r^{n-1}\dd r
= V_n\cdot\frac n2\,\Gamma\Bigl(\frac n2\Bigr)
= V_n\,\Gamma\Bigl(\frac n2 + 1\Bigr) = \pi^{n/2} = I_n .
\]
Lalu karena kedua ruasnya terbukti, kesamaannya dapat \emph{dibaca} sebagai
penguraian cangkang bagi integral Gaussnya: yakni bola
berjari-jari $r$ mengusung \omterm{def:b2:multint:domain}{luas} $nV_nr^{n-1}$, dan
bobot Gaussnya $\eu^{-r^2}$ diintegralkan atas cangkangnya.

\textbf{14.} Berlaku $s_0 = V_1 = 2$ (sebab bola-$0$-nya berupa dua titik),
$s_1 = 2V_2 = 2\pi$, $s_2 = 3V_3 = 4\pi$, dan $s_3 = 4V_4 =
2\pi^2$; lalu $\int_0^R s_{n-1}r^{n-1}\dd r = V_nR^n =
v_n(R)$: jadi luasnya adalah turunan radial volumenya.

\textbf{15.} Untuk $n = 2k$, Stirling
(\cref{thm:b2:comparison:stirling}) memberi $k! \sim
\sqrt{2\pi k}\,(k/\eu)^k$, jadi
\[
V_{2k} = \frac{\pi^k}{k!}
\sim \frac{(\pi\eu/k)^k}{\sqrt{2\pi k}}
= \frac1{\sqrt{\pi n}}\Bigl(\frac{2\pi\eu}n\Bigr)^{n/2}
\qquad (n = 2k).
\]
Sedangkan untuk $n$ yang ganjil: $V_{2k+1} = 2W_{2k+1}V_{2k} \leq 2V_{2k}$, jadi
batas peluruhan supergeometri yang sama berlaku (hingga faktor
$2$ dan pergeseran satu pada eksponennya) --- sehingga untuk setiap $q >
0$ berlaku $V_n = o(q^n)$.

\textbf{16.} Berlaku $V_2/4 = \pi/4 \approx 0.785$; $V_3/8 = \pi/6
\approx 0.524$; dan $V_{10}/2^{10} = \frac{\pi^5}{120\cdot1024}
\approx 0.0025$. Secara umum $\frac{V_n/2^n}{V_{n-2}/2^{n-2}}
= \frac{2\pi}{4n} = \frac{\pi}{2n} \to 0$: jadi nisbahnya menuju
$0$ (secara supergeometri). Sehingga bola yang terlukis di dalamnya menempati
pecahan yang lenyap: sebab volume kubusnya berpindah ke
pojoknya.

\textbf{17.} Menurut pertanyaan 1, bola dalam berjari-jari $1 -
\varepsilon$ bervolume $V_n(1-\varepsilon)^n$, jadi cangkang
luarnya mengusung pecahan $1 - (1 - \varepsilon)^n \to 1$.
Untuk $n = 100$ dan $\varepsilon = 0.01$: berlaku $(0.99)^{100} =
\eu^{100\ln0.99} \approx \eu^{-1.005} \approx 0.366$: jadi sekitar
$63\%$ bagian bolanya terletak dalam $1\%$ dari permukaannya.

\textbf{18.} Barisan $(W_n)$ turun, jadi $W_n \leq W_{n-1} \leq
W_{n-2} = \frac{n}{n-1}W_n$: sehingga lewat penjepitan, $W_{n-1}/W_n \to 1$.
Lalu mengalikannya dengan $W_nW_{n-1} = \frac\pi{2n}$: berlaku $W_n^2 \sim
\frac\pi{2n}$, yakni $W_n \sim \sqrt{\pi/(2n)}$. Adapun batas bawahnya:
$W_n^2 \geq W_nW_{n+1} = \frac{\pi}{2(n+1)}$, jadi $W_n \geq
\sqrt{\pi/(2(n+1))}$ untuk setiap $n$.

\textbf{19.} Pembilangnya: syarat $1 - x^2 \leq \eu^{-x^2}$ memberi $(1 -
x^2)^{\frac{n-1}2} \leq \eu^{-(n-1)x^2/2}$, jadi dengan $a =
\frac{n-1}2$,
\[
\int_\delta^1(1 - x^2)^{\frac{n-1}2}\dd x
\leq \int_\delta^\infty\eu^{-ax^2}\dd x
\leq \int_\delta^\infty\frac x\delta\,\eu^{-ax^2}\dd x
= \frac{\eu^{-a\delta^2}}{2a\delta}
= \frac{\eu^{-(n-1)\delta^2/2}}{(n-1)\delta}.
\]
Adapun penyebutnya: $2W_n \geq \sqrt{2\pi/(n+1)}$ menurut pertanyaan 18.
Jadi membaginya memberi batas pada displainya. Lalu untuk $\delta =
s/\sqrt{n-1}$ ia menjadi
$\sqrt{\tfrac{n+1}{2\pi(n-1)}}\;\eu^{-s^2/2}/s = O\bigl(
\eu^{-s^2/2}/s\bigr)$, secara seragam dalam $n$: jadi di luar lempeng
$\abs{x_1} \leq s/\sqrt{n-1}$ hampir tak ada volumenya, untuk
$s$ yang cukup besar --- dan berkat kesimetriannya hal yang sama berlaku bagi
setiap arah.

\textbf{20.} Kedua pernyataannya berdampingan karena keduanya memerikan
koordinat yang berbeda pada titik yang sama. Hampir setiap titik
$B_n(1)$ bernorma dekat ke $1$ (yakni pertanyaan 17: pemusatan
radial di dekat bolanya), namun tiap-tiap $n$
koordinatnya bernilai kecil, berorde $1/\sqrt n$ (yakni pertanyaan 19),
dan itu konsisten karena $n$ koordinat berukuran $1/\sqrt n$
bernorma berorde $1$. Jadi volume berdimensi tinggi memusat
di tempat semua koordinatnya berbagi anggaran \omterm{def:b2:nvs:norm}{normanya} secara merata ---
yakni di dekat bolanya, tetapi jauh dari setiap kutub sumbu koordinatnya.

\textbf{21.} Irislah $\Delta_n$ di $x_n = t \in \intcc01$: maka
irisannya adalah $\{x' \in \R^{n-1} : x_i \geq 0,\ \sum x_i \leq 1 -
t\} = (1-t)\Delta_{n-1}$, yang bervolume
$(1-t)^{n-1}\operatorname{vol}(\Delta_{n-1})$ berkat kehomogenannya.
Jadi
\[
\operatorname{vol}(\Delta_n) =
\operatorname{vol}(\Delta_{n-1})\int_0^1(1 - t)^{n-1}\dd t =
\frac{\operatorname{vol}(\Delta_{n-1})}{n}
\quad\Longrightarrow\quad
\operatorname{vol}(\Delta_n) = \frac1{n!}\,.
\]

\textbf{22.} Adapun $2^n$ oktan tandanya memotong $C_n$ menjadi $2^n$
salinan $\Delta_n$ (dengan tumpang tindih yang terabaikan pada
hiperbidang koordinatnya): jadi $\operatorname{vol}(C_n) = \frac{2^n}{n!}$. Lalu bila
$\sum\abs{x_i} \leq 1$ maka $\sum x_i^2 \leq
\bigl(\sum\abs{x_i}\bigr)^2 \leq 1$: sehingga $C_n \subseteq B_n(1)$;
dan $B_n(1) \subseteq \intcc{-1}1^n$ karena $\abs{x_i} \leq
\norm x$. Karena itu $\frac{2^n}{n!} \leq V_n \leq 2^n$ ---
yang konsisten dengan pertanyaan 15, yang menempatkan $V_n$ di antara skala
faktorial dan skala geometri.

\textbf{23.} Untuk $(x, y)$ pada cakram satuannya, irisan
$B_4(1)$ adalah cakram berjari-jari $\sqrt{1 - x^2 - y^2}$ pada
bidang-$(z, w)$, yang berluas $\pi(1 - x^2 - y^2)$. Jadi dalam koordinat
kutub:
\[
V_4 = \iint_{x^2+y^2\leq1}\pi(1 - x^2 - y^2)\,\dd x\,\dd y
= \pi\int_0^{2\pi}\!\!\int_0^1(1 - \rho^2)\rho\,
\dd\rho\,\dd\alpha
= \pi\cdot2\pi\cdot\frac14 = \frac{\pi^2}2 ,
\]
yang sepakat dengan pertanyaan 6.

\textbf{24.} Peluangnya adalah nisbah volumenya
$\dfrac{V_{20}}{2^{20}} = \dfrac{\pi^{10}}{10!\cdot2^{20}}
\approx \dfrac{0.0258}{1\,048\,576} \approx
2.5\cdot10^{-8}$. Adapun banyaknya penarikan sampai sasaran pertamanya
berorde kebalikannya, yakni sekitar $4\cdot10^7$: jadi pencuplik penolakan
yang bekerja indah bagi cakramnya (dengan $\pi/4$ sasarannya) menjadi
tak berguna pada dimensi $20$ --- yakni kutukan dimensi dalam
satu baris.

\textbf{25.} Rute pertamanya (Bagian I--II) memakai: pengirisan bertipe
Fubini pada integral teriterasinya, penyulihan satu peubah
pada tiap koordinatnya (yakni kehomogenannya), dan integral Wallisnya
--- jadi kalkulus satu peubah murni ditambah induksi. Sedangkan rute kedua
(Bagian III) memakai: Fubini bagi struktur hasil kali $I_n$,
penggantian peubah kutubnya lewat integral Gauss pada
\cref{ex:b2:multint:polar}, dan persamaan fungsional fungsi
$\Gamma$. Keduanya bertemu di $V_n =
\pi^{n/2}/\Gamma(\frac n2 + 1)$, dengan Stirling
(\cref{thm:b2:comparison:stirling}) mengubah rumusnya
menjadi asimtotik. Adapun jilid Tahun ke-3 membangun ulang semua ini atas
integral Lebesgue: sebab di sana Fubini dan penggantian peubahnya
menjadi teorema bagi fungsi terintegralkan yang umum, koordinat bolanya
ada pada setiap dimensi, dan \omterm{pb:b2:multint:1}{volume bola} yang sama
muncul kembali sebagai panen soal ukuran hasil kali dan soal
Stirling --- dengan kekonvergenan terdominasi menggantikan penjepitan buatan tangan kita.

\end{solution}
